\section{Restrição ao tamanho da dimensão extra}
Atualmente há muitos modelos que tentam impor restrições e descrevem as propriedades e o tamanho das dimensões extras \cite{MurataTanaka2015}. No entanto, esses modelos apresentam desafios signifcantes para verificação experimental devido ao grau de precisão necessário para detectar dimensões extras de tamanho pequeno. Como mencionado anteriormente neste trabalho, de acordo com a teoria de Kaluza-Klein, o tamanho da dimensão extra estaria próximo da escla de Planck, tornando difícil, senão impossível, de fazer medidas diretas usando aparato experimental atual. Aqui, porém, propomos um método indireto para detectar a existência dessa dimensão extra e derivar restrições acessíveis tanto no tamanho da dimensão extra quanto nos dois parâmetros que colocamos no nosso modelo: a razão do éter $\alpha_\phi$ e o parâmetro quasiperiódico $\beta$. 